\histitem{连续谱的发现}

在关于拓扑向量空间的注记中（参见 ÉHM, pp. 262–263），我们描述了 I. 弗雷德霍姆如何在 1900 年前后——继承 C. 诺伊曼、V. 沃尔泰拉与 H. 庞加莱的工作——开始考虑方程

\[
u (x) - \lambda \int_ {\mathrm{I}} \mathrm{K} (x, y) u (y) d y = f (x)\tag{1}
\]

其中未知函数为 \(u\colon I\to\mathbf{C}\) ，这里区间 \(I\subset\mathbf{R}\) 是紧的，且核 \(K\) 在 \(I\times I\) 上连续。他巧妙地使用“无穷行列式”——其基础是庞加莱与 H. 冯·科赫的思想——得到了对复变量 \(\lambda\) 的依赖为亚纯的解族，进而证明了关于解的存在的著名“择一定理”（参见 III, p.~81, 定理 5）。但弗雷德霍姆随即把他的结果特殊化到 \(\lambda = -1\) 的情形，而没有利用方程 (1) 是一个特征值问题 \(^{(1)}\) 这一事实。

当核 K 对称时，我们还看到希尔伯特在其关于积分方程的第一篇论文 [31] 中如何认识到弗雷德霍姆问题与求二次型的“主轴”之间的亲缘关系。从核 K 的有限“截段”（离散化）的这种约化出发，他通过取极限得到公式

\[
\int_ {\mathrm{I}} \mathrm{K} (s, t) x (s) x (t) d t = \sum_ {n = 1} ^ {\infty} \frac {1}{\lambda_ {n}} \int_ {\mathrm{I}} \varphi_ {n} (s) x (s) d s,\tag{2}
\]

其中 \(\lambda_{n}\) 是核 K 的特征值，\(\varphi_{n}\) 构成相伴特征向量 \(^{(2)}\) 的规范正交系，并且只要 x 在 I 上平方可积等式就成立（参见~ÉHM, p.~264）。

希尔伯特面临的主要障碍涉及从有限和过渡到公式 (2) 的积分。他从这一公式推出了特征值的存在性——不过，他受有限维中确定特征值的经典方法的启发（参见~IV, p.~153 第 4 目），为后者给出了一个与取极限无关的变分刻画。E. 施密特在 1905 年将证明，借鉴 H. 施瓦茨的工作，如何不经过弗雷德霍姆与希尔伯特的方法所依赖的“无穷行列式”而得到 (2)（参见~ÉHM, p.~265）。

但在 1906 年，希尔伯特 [32] 已转向约化二次型

\[
\mathrm{B} (x, x) = \sum_ {p, q = 0} ^ {\infty} b _ {p q} x _ {p} x _ {q}, \qquad x = (x _ {p}) _ {p \in \mathbf {N}}
\]

这一更一般的问题，二次型定义在 \(E = \ell_{\mathbf{C}}^{2}(\mathbf{N})\) 上。他心里显然想着：通过过渡到 E 的规范正交基下的系数，问题 (1) 可化为求“无穷线性方程组”

\[
x _ {p} + \sum_ {q = 0} ^ {\infty} b _ {p q} x _ {q} = f _ {p} \quad (p \in \mathbf {N}).
\]

的解。然而他注意到，弗雷德霍姆问题所处理的是非常特殊的二次型：\(\mathrm{B}(x_{n}, x_{n})\) 趋于 \(\mathrm{B}(x, x)\) （只要 \(x_{n}\) 弱收敛于 x）；他称之为“全连续的”（vollstetig）。

希尔伯特强调它们与 E 上有界的二次型之间的本质差别，并决定着手对后者作更一般的约化。他的方法依旧是从“有限截段”

\[
(x _ {0}, \ldots , x _ {n}) \mapsto \sum_ {p, q = 0} ^ {n} b _ {p q} x _ {p} x _ {q}
\]

的约化出发，然后让 n 趋于无穷。为完成这一取极限的过程，他利用了 T. 斯蒂尔杰斯在其关于连分数的工作 [71] 中引入的积分思想。他证明每个有界二次型经过变量的一个正交变换后都可以写成

\[
\sum_ {p \in {\bf N}} \frac {1}{\lambda_ {p}} x _ {p} ^ {2} + \int_ {s \in {\bf R}} \frac {1}{s} d \sigma (s, x).\tag{3}
\]

的形式。在这个公式中，\(\lambda_{p}\) 是 B 的特征值，\(x\mapsto\sigma(s,x)\) （对固定的 \(s\in\mathbf{R}\) ）是 E 上正定的分离二次型，而 \(s\mapsto\sigma(s,x)\) \(^{(3)}\) （对固定的 \(x\in\mathrm{E}\) ）是连续的递增函数，它在 \(-\infty\) 处趋于 0，其极限值 \(\|x\|^2\) 在 \(+\infty\) 处达到（斯蒂尔杰斯记号参见 ÉHM, p. 284）。

用 S 表示 \(\overline{R}\) 中这样的点所成的集合：不存在使 \(s \mapsto \sigma(s, x)\) 对所有 x 都保持常数的邻域；于是 (3) 中的积分自然可以在 S 上取。希尔伯特把集合 S 称为“连续谱”（Streckenspektrum），把 B 的特征值所成的集合称为“点谱”（Punktspektrum），把两者的并称为“谱”（Spektrum）。这个来自光学术语的词，曾在 1897 年被 W. 维尔丁格 [91] 用于带周期系数的微分方程的研究。

希尔伯特很快便从他的“谱”方法中为经典分析的若干问题汲取了深刻的革新，这些问题涉及常微分方程（特别是 Sturm–Liouville 型）、偏微分方程或复变量函数的研究。我们在这里不复述二十世纪头二十五年间随之而来的丰富成果，而请读者参看 E. 海林格与 O. 托普利茨的综述 [29]。不过我们仍要提到其中一些关于积分方程的工作，在它们之中可以看到抽象理论后来发展的萌芽。

对希尔伯特及其学生而言——他们忠实于线性代数中的德国传统——模型始终是主轴定理；因此所研究的总是双线性型或二次型的谱，特别是在“希尔伯特”空间 \(E = \ell_{\mathbf{C}}^{2}(\mathbf{N})\) 上。可以把 E 上任何连续双线性型 B 对应于由 \(\langle Ax, y \rangle = B(x, y)\) 刻画的 E 的连续自同态 A，这一事实为希尔伯特学派所知，尤其是在 E. 施密特、M. 弗雷歇与 F. 里斯 1907–1908 年的工作之后。但正是 F. 里斯在 1913 年一本令人赞叹的书 [59] 中指出了在此背景下把自同态置于首位的益处。

里斯给出了 \(\mathcal{L}(\mathrm{E})\) 的元素 A 的谱的现代定义 [59, p.~139]，指出它是一个有界闭集，并证明预解式 \(\lambda \mapsto (\lambda 1_{\mathrm{E}} - \mathrm{A})^{-1}\) 在 A 的谱的余集上是全纯的。尤其重要的是，他赋予 \(\mathcal{L}(\mathrm{E})\) 的代数结构以中心地位：借鉴沃尔泰拉的工作，他对对称自同态发展了泛函演算的第一个版本，使他能够简单地重新得到希尔伯特与施密特的约化结果。若 A 是 \(\mathcal{L}(\mathrm{E})\) 中的对称元素，里斯说明如何定义 \(f(\mathrm{A})\) ——对每个函数 \(f\) ，它在 \(\mathbf{R}\) 上（下或上）半连续、取值于 \(\mathbf{R}\) ——并指出 \(f \mapsto f(\mathrm{A})\) 就是我们今天所说的代数态射 [59, pp.~129--130]。把这一思想应用于区间 \([\xi, +\infty[\) 的特征函数，就对每个 \(\xi\) 得到 E 的一个对称自同态 \(\mathrm{A}_{\xi}\) ；借助斯蒂尔杰斯积分，里斯随后证明了希尔伯特关系式 (3) 的一个版本：

\[
\langle \mathrm{A} x, y \rangle = \int_ {\mathbf {R}} \xi d \langle \mathrm{A} _ {\xi} x, y \rangle
\]

对 \(x, y \in \mathrm{E}\) 成立，他甚至把这个等式写成

\[
\mathrm{A} = \int_ {\mathbf {R}} \xi d \mathrm{A} _ {\xi}.
\]

于是 A 的谱是这样的值 \(\xi_{0}\) 所成的集合：\(A_{\xi}\) 在 \(\xi_{0}\) 的任何邻域内都不保持常数；而点谱是 \(\xi \mapsto A_{\xi}\) 的不连续点所成的集合。

在以理论的应用作结之前不久，里斯指出 [59, p.~146]，当自同态 A 来自希尔伯特的某个“全连续”双线性型时，其结果具有特别简单的形式：谱约化为特征值所成之集（可能还须加上 0），特征值排成一个趋于 0 的序列，且每个非零特征值的重数是有限的（III, p.~90, 命题 5）。

